% =============================================================================
% PQ-Shield journal: replacement Abstract, Results, Discussion, Recommendations,
% and Conclusion, written from the clean 2026-10-02 re-run.
% Every number below is copied from paper/key_numbers.md (regenerate with
% `python -m analysis.paper_figures`). Figure/table files are in paper/figures
% and paper/tables; labels match the \label{} used there or suggested here.
% =============================================================================

% -----------------------------------------------------------------------------
% ABSTRACT (replaces the Review-2 abstract)
% -----------------------------------------------------------------------------
\begin{abstract}
AI inference APIs are protected almost entirely by classical key exchange and signatures, which a
cryptographically relevant quantum computer (CRQC) would break retroactively: traffic harvested
today can be decrypted later. We present PQ-Shield, a benchmarking framework that wraps a real
FastAPI inference service in six configurations --- an unprotected control, legacy RSA key transport,
production-style X25519, ML-KEM-768, the deployed X25519MLKEM768 hybrid, and ML-KEM-768 with
ML-DSA-65 signatures --- and measures them under concurrency (10--1,000 connections), four payload
shapes, emulated network paths, session reuse, and active and passive attacks, with streaming
responses from a real Llama-3.2-3B model. Against the realistic X25519 baseline, post-quantum key
establishment adds no measurable latency at low load (+3.3--7.8\% versus control for every
non-RSA configuration, including X25519 itself), costs about 5\,ms on a 5\,Mbit/s link, and its
residual overhead falls to 0--2\% when one key exchange serves 100 requests. Per-handshake RSA-2048 key
generation, by contrast, consumes 4.8 CPU cores at 10 connections and nearly doubles latency. For
token-streaming responses we show that per-chunk and hash-chained signing are the endpoints of one
design space --- a hash chain checkpointed every $k$ chunks --- and measure its trade-off between
signature bytes and \emph{exposure}, the tokens shown before verification: with ML-DSA-65 a 200-token
response costs 135,669\,B of signatures at $k{=}1$ and 3,309\,B at $k{=}\infty$, while checkpoints
move detection of dropped or reordered chunks from the end of the stream to the next checkpoint.
Streaming also compounds harvest-now-decrypt-later exposure: under classical key exchange a future
CRQC could decrypt 11.4 bytes per generated token ($R^2{=}0.9997$); under ML-KEM, none. All
post-quantum primitives are validated against NIST ACVP known-answer tests (75/75).
\end{abstract}

% -----------------------------------------------------------------------------
% §VI  RESULTS (replaces §VI A--E)
% -----------------------------------------------------------------------------
\section{Results}\label{sec:results}

All results come from one host (Apple M3, 8 cores, 16\,GB) on a single day, with every configuration
measured in isolation. Latency is end-to-end (\texttt{total\_ms}: handshake round trip plus protected
request), the only quantity directly comparable to the control's single request. Unless stated
otherwise, each cell has five repetitions.

\subsection{Cost of protection under concurrency}\label{sec:res-concurrency}
Table~\ref{tab:concurrency} and Fig.~\ref{fig:concurrency} summarize the sweep. At 10 connections
every configuration except Classical-RSA sits within a few milliseconds of the control (median
91\,ms): Classical-ECDHE +4.9\%, Hybrid +4.4\%, Full PQC +3.3\%, and Hybrid-KEX +7.8\% (all
$p<10^{-3}$, Mann--Whitney $U$). Replacing X25519 with ML-KEM-768 therefore costs nothing measurable;
the remaining few percent are the extra handshake round trip, common to all protected
configurations. Classical-RSA is the outlier at +86.5\% (median 169\,ms), and the latency
decomposition (Fig.~\ref{fig:decomposition}) places almost all of it in the handshake, where the
server generates a fresh RSA-2048 key pair.

At 100 connections the per-repetition medians of the non-RSA protected configurations are
\emph{bimodal} --- each repetition settles near either $\approx$1,000\,ms or $\approx$1,550\,ms
(e.g.\ Classical-ECDHE 1,014--1,701\,ms, Full PQC 988--1,604\,ms) --- a server-queueing regime effect
that dwarfs any cryptographic difference. We therefore do not rank these configurations at 100
connections; only Classical-RSA (+152.6\%) is separated. At 1,000 connections the four non-RSA
protected configurations are indistinguishable (medians 19.7--20.5\,s, +108.6--117.6\% versus
control), Classical-RSA reaches 30.7\,s (+225.9\%), and failures stay at 1.2--1.6\%. That the
non-RSA configurations cost about twice the control regardless of algorithm points to the second
HTTP request per transaction, not to the cryptography; Section~\ref{sec:res-amortization} confirms
that removing the per-request handshake removes the cost.

\paragraph*{Revisiting the RSA availability failure} In the Review-2 draft, Classical-RSA failed
100\% of requests at 1,000 connections on a single-core host. On the 8-core host it failed only
1.5\%, while consuming 3.5--4.8 cores of server CPU (Table~\ref{tab:resources}) versus
0.6--0.9 for every other protected configuration. Per-handshake RSA key generation is therefore
CPU-bound: it collapses when cores run out and survives when they do not, at several times the
compute of any alternative. The primitive benchmark explains why: RSA-2048 key generation runs at
16 operations/s (median 62.7\,ms over 200 samples) against 98,357/s for ML-KEM-768 --- a
6,167$\times$ gap (Table~\ref{tab:primitives}).

\subsection{Server resources}\label{sec:res-resources}
Except for Classical-RSA, CPU and memory are flat across configurations: at 10 connections every
non-RSA configuration, including the control, uses 0.92--0.94 cores and 231--241\,MB RSS
(Table~\ref{tab:resources}). The 1,088--1,120\,B ML-KEM ciphertexts and 3,309\,B ML-DSA signatures
do not add a measurable memory footprint at this scale.

\subsection{Payload shape}\label{sec:res-payload}
Across the four payload profiles (163\,B--8\,kB) the cryptographic overhead is roughly constant in
absolute terms (Fig.~\ref{fig:payload}, Table~\ref{tab:payload}): +4.0--9.4\,ms for every
non-RSA configuration and +71.6--110.2\,ms for Classical-RSA. In relative terms it ranges from
+4.5\% for the 89\,ms RandomForest to up to +240\% for millisecond-scale synthetic profiles. The cost of
protection is thus an additive term independent of payload size; whether it matters depends on the
inference time it is added to. Configurations using X25519 (Classical-ECDHE, Hybrid-KEX) are
consistently 1.6--3.5\,ms slower than those using ML-KEM alone, matching the primitive costs:
server-side handshake work is 217.9\,\textmu s for X25519 keygen plus exchange against 23.3\,\textmu s
for ML-KEM keygen plus decapsulation in our implementation.

\subsection{Network conditions}\label{sec:res-network}
On emulated paths (Fig.~\ref{fig:network}) the handshake leg isolates the cost of key-exchange size.
Relative to Classical-ECDHE, the ML-KEM configurations add $+4.6\pm0.5$\,ms (Hybrid),
$+4.7\pm0.3$\,ms (Full PQC), and $+5.7\pm1.3$\,ms (Hybrid-KEX) on the mobile path (100\,ms RTT,
5\,Mbit/s), and at most $1.0$\,ms on the metro and WAN paths. End-to-end, these differences are smaller
than the between-repetition variability of server queueing (SD 15--26\,ms on WAN and mobile), so
they are not resolvable at the application level. Classical-RSA's handshake is 61--81\,ms slower
than Classical-ECDHE's on every path because of key generation, not bytes. Because the proxy does
not model TCP slow start or loss, these are lower bounds (Section~\ref{sec:netem}).

\subsection{Session reuse}\label{sec:res-amortization}
Fig.~\ref{fig:amortization} reports overhead versus control as one key exchange serves $N$
requests. Classical-RSA's +101.2\% at $N{=}1$ disappears at $N{=}10$ ($-1.1\pm1.0$\%). For the
other configurations the overhead falls from +4.0--7.0\% at $N{=}1$ to $-0.2$ to $+2.1$\% at
$N{=}100$. One intermediate point (Full PQC at $N{=}10$, $+9.0\pm8.4$\%) has a single outlying
repetition. For any deployment that reuses sessions, post-quantum key establishment costs nothing
measurable.

\subsection{Streaming signature strategies}\label{sec:res-streaming}
With the real Llama-3.2-3B backend ($\approx$40 tokens/s; Fig.~\ref{fig:streaming},
Table~\ref{tab:streaming}), buffer-and-sign delays the first token to the end of generation: about
5\,s for a 200-token response in every configuration (4,927--5,311\,ms). Per-chunk and hash-chain
signing both deliver the first token in 125--133\,ms, a 39--42$\times$ improvement that is
independent of the key exchange. The strategies differ in signature bytes: for Full PQC a
200-token response carries 132,360\,B of ML-DSA-65 signatures under per-chunk signing versus
3,309\,B under hash-chain, a 40$\times$ difference; with ECDSA the corresponding figures are
2,840\,B and 71\,B. The cold-start artifact in the Review-2 data (13\,s for the first measured cell)
is gone after adding a discarded warm-up generation: Classical-RSA's buffer-and-sign TTFT is now
4,927\,ms, in line with the other configurations.

\subsection{Bounded exposure with checkpointed hash chains}\label{sec:res-checkpoint}
Fig.~\ref{fig:checkpoint} and Table~\ref{tab:checkpoint} sweep the checkpoint interval $k$ for
200-token responses (5-token chunks). Signature bytes follow the analytic model
$(\lfloor n/k\rfloor+1)\cdot|\sigma|$ exactly: for Full PQC 135,669\,B (41 signatures) at $k{=}1$,
29,781\,B at $k{=}5$, 9,927\,B at $k{=}20$, and 3,309\,B at $k{=}\infty$; for Hybrid 2,910\,B to
70\,B. Exposure grows as bytes fall: the client displays at most 5, 25, 100, and 200 tokens
before a signature covers them, with maximum verification lags of 0\,ms, about 0.5\,s, 2.4\,s, and 4.9\,s.
The attack experiments show the security effect directly. Without checkpoints, hash-chain detects
every mid-stream drop and reorder (100\%) but only at the end of the stream (0\% mid-stream). With
checkpoints every attack is caught mid-stream (100\%), after 50--51\% of the stream at $k\le2$
(the attack position itself), 59--60\% at $k{=}5$, 72\% at $k{=}10$, and 97--98\% at $k{=}20$.
Detection depends only on checkpoint placement, so Hybrid and Full PQC give identical curves.

\subsection{Harvest-now-decrypt-later exposure}\label{sec:res-hndl}
A passive adversary captures 209--433\,B per classical transaction and 1,265--4,503\,B per
post-quantum transaction (Fig.~\ref{fig:wirebytes}). Only the classical captures --- RSA-OAEP and
X25519 alike --- become decryptable once a CRQC exists. Under streaming
(Fig.~\ref{fig:streaminghndl}) the decryptable volume grows linearly with response length for both
classical key exchanges, at 11.4 bytes per generated token ($R^2{=}0.9997$, 50--585 tokens, the model
ending longer requests early), while all three ML-KEM configurations remain at zero decryptable bytes
at every length. Every streamed token under classical key exchange adds to a future adversary's
take; under ML-KEM, none does.

\subsection{Active tampering}\label{sec:res-mitm}
Every tampered response was rejected in every configuration (Table~\ref{tab:threats}). A flipped
ciphertext bit is caught by the AES-GCM tag in 0.007--0.008\,ms (median), before the signature is
examined; a flipped signature bit is caught by verification in 0.059\,ms (ML-DSA-65) to
0.083\,ms (ECDSA P-256). Per-chunk signing with the chunk index bound into the signed bytes detects
drops and reorders at the affected chunk (100\% mid-stream).

% -----------------------------------------------------------------------------
% §VIII  Hypotheses (replaces the H1--H4 paragraph)
% -----------------------------------------------------------------------------
\paragraph*{Hypotheses} \textbf{H1} (bounded PQC overhead) is supported: against Classical-ECDHE,
moving to ML-KEM-768 adds no measurable latency at 10 or 1,000 connections and $\le$5.7\,ms of
handshake time on the slowest emulated link; at 100 connections queueing bimodality prevents
ranking. \textbf{H2} (Hybrid captures Full PQC's HNDL benefit at lower cost) is supported on HNDL
--- both have zero decryptable bytes --- but the ``lower cost'' part is not: Full PQC is no slower
than Hybrid in any non-streaming experiment, and it costs more only in per-chunk signature bytes.
\textbf{H3} (categorically different decryptability) is demonstrated: classical exposure grows
linearly with streamed tokens, post-quantum exposure stays at zero. \textbf{H4} (ML-DSA-65
verification not slower than ECDSA P-256) is confirmed: 28,202 versus 21,602 verifications/s
(1.31$\times$).

% -----------------------------------------------------------------------------
% §IX  Recommendations (replaces §IX A--B and Table IV)
% -----------------------------------------------------------------------------
\section{Recommendations}
\paragraph*{Key establishment} Migrate from X25519 to a post-quantum key exchange now: it closes the
harvest-now-decrypt-later exposure at no measurable latency cost, a few milliseconds on slow
links, and nothing once sessions are reused. Where hedging against an ML-KEM weakness matters,
the deployed X25519MLKEM768 hybrid costs about 3\,ms more than ML-KEM alone at 10 connections in our
implementation. Do not generate RSA keys per handshake: it costs 4--5 CPU cores at modest load and
fails outright when cores run out.

\paragraph*{Signatures} ML-DSA-65 verification is faster than ECDSA P-256, so Full PQC costs no
latency for single-shot responses. Its 3,309\,B signature matters only where many signatures are
sent, which in practice means streaming.

\paragraph*{Streaming} Use hash-chained signing with checkpoints and choose $k$ from the exposure
the application can tolerate: $k=\infty$ (one signature, whole response unverified until it ends)
for low-risk display, small $k$ where a user must not act on unverified text. With ML-DSA-65,
$k{=}5$ bounds exposure to 25 tokens for 29,781\,B per 200-token response, against 132,360\,B for
per-chunk signing. Buffer-and-sign is justified only when chunk-count and timing metadata must be
hidden~\cite{weiss2024keylogging}, at a 39--42$\times$ time-to-first-token penalty.

\begin{table}[t]
\centering
\caption{Recommendation summary.}
\label{tab:recommendations}
\setlength{\tabcolsep}{3pt}
\begin{tabular}{@{}lll@{}}
\toprule
Scenario & Key exchange & Signing \\
\midrule
Non-streaming inference API & ML-KEM-768 or X25519MLKEM768 & ECDSA or ML-DSA-65 \\
Streaming, display-only & same & hash-chain, $k=\infty$ \\
Streaming, actionable output & same & hash-chain, small $k$ \\
Long-term non-repudiation & same & ML-DSA-65 (Full PQC) \\
Metadata-sensitive streaming & same & buffer-and-sign \\
Legacy RSA key transport & migrate away & -- \\
\bottomrule
\end{tabular}
\end{table}

% -----------------------------------------------------------------------------
% §XI  Conclusion (replaces the first two paragraphs)
% -----------------------------------------------------------------------------
\section{Conclusion}
Measured on a real inference service rather than a generic web server, post-quantum migration is
cheap where it is most urgent. Replacing X25519 with ML-KEM-768 closes harvest-now-decrypt-later
exposure at no measurable latency, a few milliseconds of handshake time on slow links, and no cost at
all once sessions are reused. The expensive cases are elsewhere: per-handshake RSA key generation,
which is CPU-bound and fails when cores run out, and per-chunk post-quantum signatures on streamed
responses, which multiply signature bytes by the number of chunks. For streaming, checkpointed hash
chains turn that cost into an explicit dial between signature bytes and the amount of text a client
shows before verifying it.
